\section{Attackers}
\label{sec:attackers}

% What this part is meant to establish and what is measured. What was found and why it matters are written in
% Section 7 after the results.
The attackers are meant to show what a party that holds candidate texts and knows the processing can do with what a
form leaves; what is measured is the accuracy of each attacker, and of the strongest one, on each form
(Section~\ref{sec:protocol}). Each attacker compares what a form leaves of a log with each candidate text and chooses
the candidate that fits best. Most of them compare sequences of segment values, which roughly correspond to word
lengths (\S\ref{sec:attackers:segments}), through a common alignment (\S\ref{sec:attackers:alignment}). All attackers
know the processing, and no attacker uses a column that its form does not release.

\subsection{Segments}
\label{sec:attackers:segments}

\paragraph{Texts.}
In a text, a word is a run of letters and digits, and its segment value is the number of letters and digits in it;
other characters separate words.

\paragraph{Logs.}
In a log, segments are counted from the key categories alone. A letter or digit press adds 1, a BackSpace press
subtracts 1, and a space, punctuation or Enter press closes the segment. If the next press after a Control or Meta
press, other than a modifier, is a letter or digit, it is a shortcut and is neither counted nor a separator; whether
Control is still held cannot be read, because the corpus gives modifier presses an action time of zero. Navigation,
other modifiers, mouse, Delete and other presses are neither counted nor separators, and a trailing open segment is
kept if it contains a counted press. Deletion of a selection, paste and replacement cannot be seen in the categories
and are counted by this rule as they are, so the value of a segment can differ from the length of the word it became.

\paragraph{Error of the rule.}
The checks after acquisition compared the value of each segment with the number of its counted characters that
survive in the final draft, obtained by replaying the unprocessed log. Of the 2,381,444 segments with at least one
counted press, the two values agree for 77.1\%; the share of agreeing segments per writer has a median of 0.80 (10th
percentile 0.62, 90th percentile 0.92). Empty segments, which arise for example after a comma followed by a space,
agree by construction and are not included in these shares. The error is reported with the tables, and the rule was
not changed; cases that the rule misses are described in Appendix~\ref{app:checks}.

\subsection{Alignment}
\label{sec:attackers:alignment}

Attackers A1, A1', A2 and A3 align a sequence of record segments with a sequence of text segments by dynamic
programming. A group maps $p$ record segments to $q$ text segments, with $(p, q)$ one of $(1,1)$, $(2,1)$, $(1,2)$,
$(3,1)$ and $(1,3)$, which absorbs a separator that was typed but not kept or kept but not typed. The cost of a group
is
\[
  \Bigl|\sum \text{record values} - \sum \text{text values}\Bigr| + \lambda\,(p - 1 + q - 1),
\]
with $\lambda = 8$, a value fixed in earlier prototypes and not tuned on these data; $\lambda = 4$ and $16$ are reported
as a secondary analysis. Among paths of equal cost, the group listed first is chosen, and where the end can be chosen, the earlier position in the
text. In full-text matching both
sequences are used to their ends; in window matching the whole window is used and the text may end early. Under
suppression (R7), dropping a separator press merges record segments, so the groups are extended to $(p', 1)$ and
$(1, q)$ with $p', q$ up to $\lceil 3/(1-p) \rceil$, that is 4 for $p = 0.1$ and $0.25$ and 6 for $p = 0.5$; the
cost is the same. This alignment ignores the binomial variation of suppression and is not a likelihood. A candidate
with no path has infinite cost and is never chosen; if every candidate has infinite cost, all candidates tie. The
numbers of such cases are reported.

\subsection{Attackers}
\label{sec:attackers:list}

Table~\ref{tab:attackers} lists the attackers and the forms on which each is run. Attackers A2-b and A3 are trained on
the unprocessed logs and texts of all writers of the two prompts not being evaluated, with the evaluated form applied to
the training logs. The strongest attacker of a form is the best of the adaptive attackers (A0 to A3, including A1$'$ and A2-b) on that
form
(\S\ref{sec:protocol:quantities}).

\begin{table*}
  \caption{Attackers. A0 to A3 are within the threat model and adapt to the processing; A4 is a reference outside the
  threat model.}
  \label{tab:attackers}
  \small
  \begin{tabular}{>{\raggedright\arraybackslash}p{0.16\linewidth}>{\raggedright\arraybackslash}p{0.17\linewidth}>{\raggedright\arraybackslash}p{0.58\linewidth}}
    \toprule
    Attacker & Forms & What it compares \\
    \midrule
    A0 Total length & All & Letter and digit presses minus BackSpace presses in the log (R3-u: ``char'' presses
      minus BackSpace; R7: the remaining presses scaled by $1/(1-p)$), against the number of letters and digits in the
      candidate (R3-u: all characters); in R4, the sum of the released values against the sum of $\lfloor v/w
      \rfloor$ over the candidate's segments \\
    A1 Final-draft lengths & R0, R1, R2, R5, R6 & The final draft is rebuilt from the text change and the position (R2:
      the categories and the position), as in the winning solution of the competition~\cite{inubushi2024first}, and
      its word lengths are aligned with the candidate's \\
    A1$'$ Editor replay & R3, R3-t, R7 & The category column is replayed as an editor, and the word lengths of the
      replayed draft are aligned with the candidate's (R7: candidate lengths scaled by $1-p$) \\
    A2 Segment alignment & R2, R3, R3-t, R4, R7 & The record segments are aligned with the candidate's segments (R4:
      both as $\lfloor v/w \rfloor$; R7: candidate values scaled by $1-p$) \\
    A2-b Boundary estimation & R3-u & A logistic regression predicts which ``char'' presses are separators from the
      intervals around them; the predicted segments are aligned as in A2 \\
    A3 Timing & R0, R1, R2, R3, R3-t, R3-u, R5, R6, R7 & The likelihood of the press intervals under the character
      pairs of the candidate, on the A2 alignment; in R3-u combined with a length term; in R7 the duration of each
      segment \\
    A4 KeyTAR (reference) & R0, R1, R2, R3, R3-t & The public model reconstructs text from the intervals of letter, digit
      and space presses; the candidate nearest to the reconstruction is chosen \\
    \bottomrule
  \end{tabular}
\end{table*}

\paragraph{A0: total length.}
A0 is the baseline that every other attacker is compared with: it chooses the candidate whose length is closest to
the number of characters that the log leaves.

\paragraph{A1 and A1$'$: the final draft.}
A1 rebuilds the final draft from the columns that R0, R1, R2, R5 and R6 release. A1$'$ uses only the category column:
a character press inserts at the cursor, BackSpace deletes to the left and Delete to the right, ArrowLeft, ArrowRight,
Home and End move the cursor, with lines separated by Enter, and other navigation and mouse presses put the cursor at
the end, since their effect cannot be known. After a Control press, Home and End move to the start and the end of the
document and the arrows are treated as unknown; shortcuts are ignored. Both attackers align the word lengths of the
draft with those of each candidate over the full text and choose the candidate with the lowest cost.

\paragraph{A2 and A2-b: segments.}
A2 aligns the record segments directly with the candidate's word lengths. In R3-u the separators are merged into
``char'', so A2-b first estimates them: a logistic regression, trained on the training logs processed to R3-u,
predicts from the intervals before and after each ``char'' press, and from the neighbouring intervals, whether the
press is a space, punctuation or Enter.

\paragraph{A3: timing.}
From the training logs, A3 estimates the mean and variance of the logarithm of the press interval for each pair of
consecutive characters. A pair seen fewer than 20 times takes the values of its pair of character types (such as
letter followed by space), and a type pair seen fewer than 20 times takes the values of all intervals. For each record
segment without a BackSpace, the intervals are scored under the character pairs of the candidate word that the A2
alignment assigns to it, and the log-likelihoods are summed; the candidate with the highest sum is chosen. In R3-t the
likelihood of a rounded interval is the probability of its rounding interval. In R3-u, where segments are not visible,
A3 combines two terms: the A1$'$ replay gives the sequence of ``char'' items of the final draft, matched one to one from
the start with all characters of the candidate, and the timing term uses the intervals between adjacent final
characters that were typed consecutively; the length term is the log-likelihood of the difference between the replayed
count and the candidate's character count, estimated on the training logs. The timing term is used both as a sum and
as a mean per term, and the better of the two versions is reported as A3 on R3-u; a version with the length term alone
is reported to separate the contribution of timing. In R7, A3 scores the duration of each segment, from its first to
its last remaining press, with a regression on the logarithm of the candidate segment's character count, estimated on
training logs with the same suppression. The check on synthetic logs showed that the one-to-one correspondence in R3-u
uses in effect only the part of the log before the first replay error. Details that the registration left open, such
as how an interval is scored when the candidate word has fewer characters than the segment has intervals, were fixed
on synthetic logs before the data were opened and are documented with the code.

\paragraph{A4: KeyTAR, a reference outside the threat model.}
A4 gives the intervals between letter, digit and space presses, in chunks, to the public KeyTAR model and chooses the
candidate nearest to the reconstruction. The model is used with the default settings of its artifact, which include a
limit on generated tokens that is expected to cut many reconstructions short; we kept it as registered. KeyTAR assumes
ordinary English input without revision~\cite{qiu2026keytar}, it contains a language model, and it does not adapt to
the processing. It is therefore reported in a separate row and is not part of the strongest attacker. The details are
in Appendix~\ref{app:attackers}.

\subsection{Attackers on windows}
\label{sec:attackers:windows}

A window is a run of segments of the unprocessed log that contains $k$ non-empty segments; empty segments between them
are inside the window but not counted. Its rows run from the row after the separator that precedes its first segment
to the separator that closes its last segment, and include the mouse, navigation and modifier rows in between. Each
form takes these rows and segments them with its own rule, so every form starts from the same presses. Because the
ends of a window are word boundaries of the unprocessed log, the attacker knows both ends of the window even in forms
that hide separators (R3-u) or drop presses (R7). This favours the attacker
(\S\ref{sec:threat-model:bias}), and the difference between R3 and R3-u in windows may be smaller than it would be for
parts of logs released in practice.

On windows, A1 and A1$'$ are not run, since a part of the process does not determine the final draft; the other
attackers compare the window with a candidate window, as described in Appendix~\ref{app:attackers}.

\subsection{Attackers on the Japanese translation logs}
\label{sec:attackers:enja}

In a translation, segments are split at the Japanese comma and full stop and at line breaks, and the value of a
segment is the standard number of keystrokes needed to type its reading with an input method. In a log, a letter press
adds 1, a BackSpace press subtracts 1, and a commit that contains punctuation closes the segment. Only A0 and A2 are
run on these logs, as on KLiCKe; neither reads a string, and R0j is only a reference (Appendix~\ref{app:enja}).
